%%
%% part2-book-components/ch08-environments.tex
%%
%% Chapter 8: Scientific Environments — Theorem, lemma,
%% corollary, proposition, definition, example, remark,
%% exercise, solution, and proof.
%%

\chapter{محیط‌های علمی}
\label{chap:environments}

\section{چرا این موضوع مهم است؟}
\label{sec:environments-why}

در کتاب‌های ریاضی و فنی، محتوا به انواع مختلفی
تقسیم می‌شود: قضیه، لم، تعریف، مثال، اثبات و...
هر نوع، نقش معنایی مشخصی دارد و باید در قالب
بصری متمایز نمایش داده شود تا خواننده بتواند
به‌سرعت نوع محتوا را تشخیص دهد.

\lr{MatinBook} ده نوع محیط علمی را با دو خانواده‌ی
بصری متمایز ارائه می‌دهد. در این فصل، با این
محیط‌ها آشنا می‌شوید و یاد می‌گیرید که چگونه
از آن‌ها استفاده کنید.

\mnote{استفاده‌ی درست از محیط‌های علمی، یکی از عوامل مهم در کیفیت کتاب‌های ریاضی و فنی است.}

\section{دو خانواده‌ی جعبه}
\label{sec:environments-families}

محیط‌های علمی \lr{MatinBook} به دو خانواده
تقسیم می‌شوند:

\begin{enumerate}
    \item \textbf{خانواده‌ی \lr{matinbox}:}
    با یک نوار رنگی توپر در بالای جعبه. این
    خانواده برای محیط‌های نظری و رسمی استفاده
    می‌شود.
    \item \textbf{خانواده‌ی \lr{matin-outline}:}
    با یک خط عمودی در سمت راست و دو خط افقی
    زیر عنوان و زیر محتوا. این خانواده برای
    محیط‌های توضیحی و غیررسمی استفاده می‌شود.
\end{enumerate}

\mnote{تفکیک دو خانواده، بر اساس نوع محتوا است، نه سلیقه. محیط‌های نظری «رسمی» هستند، محیط‌های توضیحی «غیررسمی».}

\subsection{خانواده‌ی \lr{matinbox}}
\label{sec:environments-matinbox}

خانواده‌ی \lr{matinbox} شامل هفت محیط است:

\begin{itemize}
    \item \env{theorem} — قضیه.
    \item \env{lemma} — لم.
    \item \env{corollary} — نتیجه.
    \item \env{proposition} — گزاره.
    \item \env{definition} — تعریف.
    \item \env{remark} — نکته.
    \item \env{exercise} — تمرین.
\end{itemize}

\subsection{خانواده‌ی \lr{matin-outline}}
\label{sec:environments-outline}

خانواده‌ی \lr{matin-outline} شامل سه محیط است:

\begin{itemize}
    \item \env{example} — مثال.
    \item \env{proof} — اثبات.
    \item \env{solution} — راه‌حل.
\end{itemize}

\mnote{خانواده‌ی \lr{matin-outline} در نسخه‌ی \lr{1.2} بازطراحی شده و خط عمودی آن همیشه در سمت راست است.}

\section{محیط‌های خانواده‌ی \lr{matinbox}}
\label{sec:environments-matinbox-detail}

\subsection{قضیه (\lr{theorem})}
\label{sec:environments-theorem}

قضیه، یک حکم ریاضی است که اثبات شده است.
برای ایجاد قضیه:

\begin{latin}
\begin{minted}{latex}
\begin{theorem}[|\pc{قضیه‌ی فیثاغورث}|]
    |\pc{در هر مثلث قائم‌الزاویه، مربع وتر برابر است با}|%
    |\pc{مجموع مربعات دو ضلع دیگر:}|%
    \[
        a^{2} + b^{2} = c^{2}
    \]
    \label{thm:pythagoras}
\end{theorem}
\end{minted}
\end{latin}

خروجی:

\begin{theorem}[قضیه‌ی فیثاغورث]
    در هر مثلث قائم‌الزاویه، مربع وتر برابر است با
    مجموع مربعات دو ضلع دیگر:
    \[
        a^{2} + b^{2} = c^{2}
    \]
    \label{thm:pythagoras-ch8}
\end{theorem}

\mnote{عنوان داخل \lr{[]} اختیاری است. اگر آن را ندهید، فقط «قضیه ۱.۱» نمایش داده می‌شود.}

\subsection{لم (\lr{lemma})}
\label{sec:environments-lemma}

لم، یک قضیه‌ی کمکی است که برای اثبات قضیه‌های
بزرگ‌تر استفاده می‌شود:

\begin{latin}
\begin{minted}{latex}
\begin{lemma}[|\pc{لم اقلیدس}|]
    |\pc{اگر}|\lr{$p$}|\pc{عددی اول باشد و}|\lr{$p \mid ab$}|\pc{، آنگاه}|%
    |\lr{$p \mid a$}|\pc{یا}|\lr{$p \mid b$}|\pc{.}|%
    \label{lem:euclid}
\end{lemma}
\end{minted}
\end{latin}

خروجی:

\begin{lemma}[لم اقلیدس]
    اگر \lr{$p$} عددی اول باشد و \lr{$p \mid ab$}،
    آنگاه \lr{$p \mid a$} یا \lr{$p \mid b$}.
    \label{lem:euclid-ch8}
\end{lemma}

\subsection{نتیجه (\lr{corollary})}
\label{sec:environments-corollary}

نتیجه، یک حکم است که بلافاصله از یک قضیه یا
لم نتیجه می‌شود:

\begin{latin}
\begin{minted}{latex}
\begin{corollary}
    |\pc{هر دنباله‌ی یکنوا و کراندار همگراست.}|%
    \label{cor:monotone}
\end{corollary}
\end{minted}
\end{latin}

خروجی:

\begin{corollary}
    هر دنباله‌ی یکنوا و کراندار همگراست.
    \label{cor:monotone-ch8}
\end{corollary}

\subsection{گزاره (\lr{proposition})}
\label{sec:environments-proposition}

گزاره، یک حکم است که ارزش اثبات دارد ولی
لزوماً به اندازه‌ی قضیه مهم نیست:

\begin{latin}
\begin{minted}{latex}
\begin{proposition}
    |\pc{مجموعه‌ی اعداد گویا}|\lr{$\Q$}|\pc{در}|\lr{$\R$}|\pc{چگال است.}|%
    \label{prop:q-dense}
\end{proposition}
\end{minted}
\end{latin}

خروجی:

\begin{proposition}
    مجموعه‌ی اعداد گویا \lr{$\Q$} در \lr{$\R$}
    چگال است.
    \label{prop:q-dense-ch8}
\end{proposition}

\subsection{تعریف (\lr{definition})}
\label{sec:environments-definition}

تعریف، یک مفهوم جدید را معرفی می‌کند:

\begin{latin}
\begin{minted}{latex}
\begin{definition}[|\pc{گراف}|]
    |\pc{گراف}|\lr{$G = (V, E)$}|\pc{مجموعه‌ای از رئوس}|\lr{$V$}|%
    |\pc{و یال‌های}|\lr{$E$}|\pc{است.}|%
    \label{def:graph}
\end{definition}
\end{minted}
\end{latin}

خروجی:

\begin{definition}[گراف]
    گراف \lr{$G = (V, E)$} مجموعه‌ای از رئوس
    \lr{$V$} و یال‌های \lr{$E$} است.
    \label{def:graph-ch8}
\end{definition}

\mnote{تعریف‌ها معمولاً در فصل‌های اول کتاب می‌آیند و پایه‌ی مفاهیم بعدی هستند.}

\subsection{نکته (\lr{remark})}
\label{sec:environments-remark}

نکته، یک توضیح مکمل یا هشدار است:

\begin{latin}
\begin{minted}{latex}
\begin{remark}
    |\pc{تقسیم بر صفر تعریف نشده است! هنگام ساده‌سازی}|%
    |\pc{عبارات کسری، همواره فرض کنید مخرج غیرصفر است.}|%
    \label{rem:division-zero}
\end{remark}
\end{minted}
\end{latin}

خروجی:

\begin{remark}
    تقسیم بر صفر تعریف نشده است! هنگام ساده‌سازی
    عبارات کسری، همواره فرض کنید مخرج غیرصفر است.
    \label{rem:division-zero-ch8}
\end{remark}

\mnote{استفاده‌ی سنجیده از نکته ضروری است. اگر بیش از حد استفاده شود، اثر هشداردهندگی آن از بین می‌رود.}

\subsection{تمرین (\lr{exercise})}
\label{sec:environments-exercise}

تمرین، یک مسئله برای حل کردن است:

\begin{latin}
\begin{minted}{latex}
\begin{exercise}
    |\pc{ثابت کنید مجموع دو عدد زوج، زوج است.}|%
    \label{exc:sum-even}
\end{exercise}
\end{minted}
\end{latin}

خروجی:

\begin{exercise}
    ثابت کنید مجموع دو عدد زوج، زوج است.
    \label{exc:sum-even-ch8}
\end{exercise}

\section{محیط‌های خانواده‌ی \lr{matin-outline}}
\label{sec:environments-outline-detail}

\subsection{مثال (\lr{example})}
\label{sec:environments-example}

مثال، یک نمونه‌ی عملی از یک مفهوم است:

\begin{latin}
\begin{minted}{latex}
\begin{example}[|\pc{محاسبه‌ی انتگرال}|]
    |\pc{مساحت زیر نمودار}|\lr{$f(x) = x^{2}$}|\pc{از}|\lr{$x = 0$}|%
    |\pc{تا}|\lr{$x = 2$}|\pc{:}|%
    \[
        \int_{0}^{2} x^{2} \, dx = \left[ \frac{x^{3}}{3} \right]_{0}^{2}
        = \frac{8}{3} \approx 2.67
    \]
    \label{ex:integral-x2}
\end{example}
\end{minted}
\end{latin}

خروجی:

\begin{example}[محاسبه‌ی انتگرال]
    مساحت زیر نمودار \lr{$f(x) = x^{2}$} از
    \lr{$x = 0$} تا \lr{$x = 2$}:
    \[
        \int_{0}^{2} x^{2} \, dx = \left[ \frac{x^{3}}{3} \right]_{0}^{2}
        = \frac{8}{3} \approx 2.67
    \]
    \label{ex:integral-x2-ch8}
\end{example}

\subsection{اثبات (\lr{proof})}
\label{sec:environments-proof}

اثبات، استدلال ریاضی برای یک حکم است:

\begin{latin}
\begin{minted}{latex}
\begin{proof}
    |\pc{از برهان خلف استفاده می‌کنیم. فرض کنید}|\lr{$\sqrt{2}$}|%
    |\pc{گویا باشد:}|\lr{$\sqrt{2} = \frac{a}{b}$}|\pc{که}|\lr{$\gcd(a,b) = 1$}|\pc{.}|%

    \[
    2 = \frac{a^{2}}{b^{2}} \implies a^{2} = 2b^{2}
    \]

    |\pc{بنابراین}|\lr{$a$}|\pc{زوج است. پس}|\lr{$b$}|\pc{نیز زوج است که}|%
    |\pc{با فرض}|\lr{$\gcd(a,b) = 1$}|\pc{تناقض دارد.}|%
\end{proof}
\end{minted}
\end{latin}

خروجی:

\begin{proof}
    از برهان خلف استفاده می‌کنیم. فرض کنید \lr{$\sqrt{2}$}
    گویا باشد: \lr{$\sqrt{2} = \frac{a}{b}$} که
    \lr{$\gcd(a,b) = 1$}.

    \[
    2 = \frac{a^{2}}{b^{2}} \implies a^{2} = 2b^{2}
    \]

    بنابراین \lr{$a$} زوج است. پس \lr{$b$} نیز
    زوج است که با فرض \lr{$\gcd(a,b) = 1$}
    تناقض دارد.
\end{proof}

\mnote{محیط اثبات به‌طور خودکار علامت پایان (□) را در انتهای جعبه اضافه می‌کند.}

\subsection{راه‌حل (\lr{solution})}
\label{sec:environments-solution}

راه‌حل، پاسخ تشریحی یک تمرین است:

\begin{latin}
\begin{minted}{latex}
\begin{solution}
    |\pc{فرض کنید}|\lr{$a = 2m$}|\pc{و}|\lr{$b = 2n$}|\pc{.}|%
    |\pc{آنگاه}|\lr{$a + b = 2(m + n)$}|\pc{که زوج است.}|%
\end{solution}
\end{minted}
\end{latin}

خروجی:

\begin{solution}
    فرض کنید \lr{$a = 2m$} و \lr{$b = 2n$}.
    آنگاه \lr{$a + b = 2(m + n)$} که زوج است.
\end{solution}

\mnote{محیط راه‌حل و اثبات شماره‌گذاری نمی‌شوند، چون به‌طور طبیعی به تمرین یا قضیه‌ی قبلی متصل هستند.}

\section{شمارنده‌ی مشترک}
\label{sec:environments-counter}

تمام محیط‌های شماره‌دار (قضیه، لم، نتیجه، گزاره،
تعریف، مثال، نکته، تمرین) از یک \textbf{شمارنده‌ی
مشترک} استفاده می‌کنند. این یعنی شماره‌گذاری
پیوسته است:

\begin{latin}
\begin{minted}{latex}
\begin{definition}[|\pc{گراف}|]   % |\pc{شماره ۱.۱}|%
    ...
\end{definition}

\begin{theorem}[|\pc{اویلر}|]   % |\pc{شماره ۱.۲}|%
    ...
\end{theorem}

\begin{lemma}[|\pc{دست‌دادن}|]   % |\pc{شماره ۱.۳}|%
    ...
\end{lemma}

\begin{example}                % |\pc{شماره ۱.۴}|%
    ...
\end{example}
\end{minted}
\end{latin}

\mnote{این شمارنده‌گذاری مشترک، یک تصمیم آگاهانه است. جامعه‌ی \lr{LaTeX} فارسی ترجیح می‌دهد شماره‌گذاری پیوسته باشد تا خواننده راحت‌تر مفاهیم را پیدا کند.}

\section{ارجاع به محیط‌ها}
\label{sec:environments-references}

برای ارجاع به محیط‌های علمی، از دستور \cmd{cref}
استفاده می‌کنیم:

\begin{latin}
\begin{minted}{latex}
% |\pc{ارجاع به قضیه}|%
\cref{thm:pythagoras}

% |\pc{ارجاع به تعریف}|%
\cref{def:graph}

% |\pc{ارجاع به چند مرجع}|%
\cref{thm:pythagoras,def:graph}

% |\pc{ارجاع بازه‌ای}|%
\crefrange{thm:pythagoras}{def:graph}
\end{minted}
\end{latin}

\mnote{دستور \lr{\textbackslash cref} به‌طور خودکار نوع مرجع را تشخیص می‌دهد و پیشوند مناسب (قضیه، تعریف، لم) را اضافه می‌کند.}

\subsection{نام‌گذاری برچسب‌ها}
\label{sec:environments-labels}

برای برچسب‌ها، از پیشوندهای زیر استفاده کنید:

\begin{itemize}
    \item \textbf{قضیه:} \lr{thm:xxx}.
    \item \textbf{لم:} \lr{lem:xxx}.
    \item \textbf{نتیجه:} \lr{cor:xxx}.
    \item \textbf{گزاره:} \lr{prop:xxx}.
    \item \textbf{تعریف:} \lr{def:xxx}.
    \item \textbf{مثال:} \lr{ex:xxx}.
    \item \textbf{نکته:} \lr{rem:xxx}.
    \item \textbf{تمرین:} \lr{exc:xxx}.
\end{itemize}

\mnote{برچسب‌های معنادار، خواندن فایل‌های \lr{tex} را آسان‌تر می‌کنند و از خطاهای ارجاع جلوگیری می‌کنند.}

\section{جدول خلاصه}
\label{sec:environments-summary-table}

جدول \cref{tab:environments-summary} خلاصه‌ای از
ده محیط علمی \lr{MatinBook} را نشان می‌دهد.

\begin{matintable}{خلاصه‌ی محیط‌های علمی \lr{MatinBook}}{tab:environments-summary}
\begin{tblr}{
    width = \textwidth,
    colspec = {l l l X[c]},
    hlines, vlines,
    colsep = 6pt,
    rowsep = 4pt,
    row{1} = {font=\bfseries},
}
محیط & عنوان & خانواده & شماره‌دار \\
\lr{theorem} & قضیه & \lr{matinbox} & بله \\
\lr{lemma} & لم & \lr{matinbox} & بله \\
\lr{corollary} & نتیجه & \lr{matinbox} & بله \\
\lr{proposition} & گزاره & \lr{matinbox} & بله \\
\lr{definition} & تعریف & \lr{matinbox} & بله \\
\lr{remark} & نکته & \lr{matinbox} & بله \\
\lr{exercise} & تمرین & \lr{matinbox} & بله \\
\lr{example} & مثال & \lr{matin-outline} & بله \\
\lr{proof} & اثبات & \lr{matin-outline} & خیر \\
\lr{solution} & راه‌حل & \lr{matin-outline} & خیر \\
\end{tblr}
\end{matintable}

\section{خطاهای رایج}
\label{sec:environments-mistakes}

\subsection{قرار دادن \cmd{label} بیرون محیط}
\label{sec:environments-mistake-label}

اگر \cmd{label} را بیرون محیط قرار دهید، ارجاع
به آن کار نمی‌کند:

\begin{latin}
\begin{minted}{latex}
% |\pc{اشتباه}|:
\begin{theorem}
    |\pc{متن قضیه...}|%
\end{theorem}
\label{thm:wrong}

% |\pc{درست}|:
\begin{theorem}
    |\pc{متن قضیه...}|%
    \label{thm:correct}
\end{theorem}
\end{minted}
\end{latin}

\subsection{فراموش کردن \cmd{label}}
\label{sec:environments-mistake-no-label}

اگر \cmd{label} را فراموش کنید، نمی‌توانید به
آن محیط ارجاع دهید.

\subsection{استفاده از محیط نامناسب}
\label{sec:environments-mistake-wrong-env}

هر محیط، نقش معنایی مشخصی دارد. استفاده از
\env{theorem} برای یک تعریف یا \env{definition}
برای یک قضیه، ساختار کتاب را به‌هم می‌ریزد.

\section{تمرین‌ها}
\label{sec:environments-exercises}

\begin{exercise}
    یک قضیه، یک تعریف و یک مثال بنویسید و
    به هر یک ارجاع دهید.
    \label{exc:environments-basic}
\end{exercise}

\begin{exercise}
    یک لم بنویسید و یک نتیجه از آن استخراج
    کنید.
    \label{exc:environments-lemma}
\end{exercise}

\begin{exercise}
    یک اثبات برای قضیه‌ی فیثاغورث بنویسید
    و از محیط \env{proof} استفاده کنید.
    \label{exc:environments-proof}
\end{exercise}

\begin{exercise}
    یک تمرین بنویسید و راه‌حل آن را با محیط
    \env{solution} ارائه دهید.
    \label{exc:environments-solution}
\end{exercise}

\section{خلاصه}
\label{sec:environments-summary}

در این فصل، با محیط‌های علمی \lr{MatinBook}
آشنا شدیم:

\begin{itemize}
    \item \lr{MatinBook} ده محیط علمی را در دو
    خانواده‌ی بصری ارائه می‌دهد.

    \item خانواده‌ی \lr{matinbox} شامل قضیه، لم،
    نتیجه، گزاره، تعریف، نکته و تمرین است.

    \item خانواده‌ی \lr{matin-outline} شامل مثال،
    اثبات و راه‌حل است.

    \item تمام محیط‌های شماره‌دار، از یک
    شمارنده‌ی مشترک استفاده می‌کنند.

    \item اثبات و راه‌حل، شماره‌گذاری نمی‌شوند.

    \item ارجاع به محیط‌ها با \cmd{cref} انجام
    می‌شود که پیشوند مناسب را خودکار اضافه
    می‌کند.

    \item خطاهای رایج شامل قرار دادن \cmd{label}
    بیرون محیط، فراموش کردن \cmd{label}، و
    استفاده از محیط نامناسب است.
\end{itemize}

در فصل بعدی (\cref{chap:math})، با ریاضیات
در \lr{MatinBook} آشنا می‌شوید و یاد می‌گیرید
که چگونه معادلات، ماتریس‌ها و عملگرها را
بنویسید.